\chapter{Roteiro de entrevista com fornecedor de cardápio digital}
\label{ap:roteiro-fornecedor}

Este apêndice apresenta o roteiro semiestruturado aplicado ao representante de uma empresa fornecedora de cardápios digitais (Seção~\ref{sec:entrevista-fornecedor}). A entrevista tem duração estimada de 40 a 50 minutos e é gravada mediante consentimento (Apêndice~\ref{ap:termos}). Os Blocos~2 e~4 são prioritários; se o tempo for curto, os Blocos~1 e~5 são abreviados. O entrevistador pede exemplos concretos e não menciona resultados do questionário nem das demais entrevistas, para não induzir as respostas.

\noindent\textbf{Bloco 0, Caracterização}
\begin{enumerate}
	\item Conte brevemente a história da empresa e do produto: quando surgiu, que problema queria resolver, onde atua.
	\item Qual é o seu papel na empresa e há quanto tempo você está nessa função?
	\item Mais ou menos quantos estabelecimentos usam a solução hoje, e quantos em Porto Alegre ou no RS? Qual é o perfil típico (porte, tipo de casa)?
\end{enumerate}

\noindent\textbf{Bloco 1, Produto e modalidades}
\begin{enumerate}\setcounter{enumi}{3}
	\item Que modalidades vocês oferecem (QR Code no celular do cliente, tablet, totem) e qual é a mais vendida? Por quê?
	\item O cliente final só consulta o cardápio ou também faz o pedido e paga pelo sistema? Qual é a proporção entre esses usos?
	\item O cardápio é uma página web, um aplicativo ou um arquivo (PDF)? Por que essa escolha?
\end{enumerate}

\noindent\textbf{Bloco 2, Arquitetura}
\begin{enumerate}\setcounter{enumi}{6}
	\item Em linhas gerais, como a solução é construída e hospedada? É uma plataforma única que atende todos os estabelecimentos (multi-inquilino) ou cada cliente tem uma instalação própria? Como os dados de cada estabelecimento são separados?
	\item Com que outros sistemas a solução se integra (PDV, emissão de nota fiscal, estoque, impressora ou tela da cozinha, meios de pagamento, delivery)? Quais integrações dão mais trabalho?
	\item O que acontece quando a internet do estabelecimento ou do cliente cai? Existe algum funcionamento offline ou plano B?
	\item Como o estabelecimento atualiza preços, itens novos e itens esgotados? Isso aparece para o cliente em tempo real?
	\item Que nível de disponibilidade vocês garantem? Já tiveram uma queda grande? O que aprenderam com ela?
	\item O que cada estabelecimento pode personalizar (visual, fotos, categorias, idiomas) e o que é igual para todos?
\end{enumerate}

\noindent\textbf{Bloco 3, Requisitos e evolução}
\begin{enumerate}\setcounter{enumi}{12}
	\item Quais são as funcionalidades que os estabelecimentos mais pedem? E quais eles quase não usam?
	\item Vocês recebem retorno dos clientes finais? Quais são as reclamações e os elogios mais comuns, e como esse retorno chega até vocês?
	\item Já lançaram alguma funcionalidade que não deu certo? Por que não funcionou?
	\item Que dados de uso vocês oferecem ao estabelecimento (itens mais vistos, horários de pico, tempo de pedido)? Os clientes usam esses dados?
\end{enumerate}

\noindent\textbf{Bloco 4, Dados pessoais e acessibilidade}
\begin{enumerate}\setcounter{enumi}{16}
	\item O cliente final precisa se cadastrar ou informar algum dado (nome, CPF, telefone) para ver o cardápio ou pedir? Quem decide isso? Para que esses dados são usados?
	\item Na LGPD, como vocês entendem o papel de cada um: o estabelecimento é o controlador e vocês o operador, ou é diferente? Como isso aparece no contrato?
	\item A solução segue alguma diretriz de acessibilidade (por exemplo, WCAG)? Funciona com leitor de tela, permite aumentar a fonte, tem bom contraste? Já foi testada com pessoas idosas ou com deficiência?
\end{enumerate}

\noindent\textbf{Bloco 5, Mercado e adoção}
\begin{enumerate}\setcounter{enumi}{19}
	\item Que tipo de estabelecimento adota mais facilmente e que tipo resiste? Quais são as objeções mais comuns na venda?
	\item Como a pandemia mudou a procura pela solução? E depois dela?
	\item A lei de Porto Alegre de 2024, que exige cardápio físico junto com o digital, afetou vendas, cancelamentos ou o uso do produto?
	\item Como funciona o modelo de cobrança (mensalidade, taxa por pedido, equipamento)? Quando um cliente cancela, quais costumam ser os motivos?
\end{enumerate}

\noindent\textbf{Bloco 6, Encerramento}
\begin{enumerate}\setcounter{enumi}{23}
	\item Daqui a cinco anos, como você imagina o atendimento em bares e restaurantes?
	\item Tem algo importante sobre esse assunto que eu não perguntei?
	\item Vocês poderiam compartilhar algum dado agregado e anonimizado (por exemplo, proporção de pedidos pelo QR Code ou tempo médio de pedido) para o trabalho?
\end{enumerate}
